\documentclass[11pt,a4paper]{article}
\usepackage[utf8]{inputenc}
\usepackage{graphicx}
\usepackage[left=2.5cm,top=3cm,right=2.5cm,bottom=3cm,bindingoffset=0.5cm]{geometry}
\usepackage{AEDLogica, AEDEspecificacion, AEDTADs}
\usepackage{caratula}
\newlength{\miSangria}
\setlength{\miSangria}{2em}

% Asi pueden escribir nuevos comandos. 
% Este por ejemplo asegura q los nombres 
% que figuren con una tipografia diferenciada  
\newcommand{\Tipo}[1]{\mathsf{#1}}
% la sintaxis es \newcommand{\nombreDeLaMacro}[cantidadDeParametros]{Lo que va ser remplazado por el macro} 
\newcommand{\norm}[1]{\vert#1\vert}


\begin{document}

\section{Titulo del Proc ACA}
\vspace{1.0cm}
\noindent Transacciones \textbf{ES} Tuple $\tupla{
		UsuarioOrigen: \Tipo{\Z}, \\
		UsuarioDestino: \Tipo{\Z},
		TipoDeTransaccion: \Tipo{\Z},
		Millas: \Tipo{\R}, \\
		NombrePromo: \Tipo{\str},
		FactorPromo: \Tipo{\R}
	}$

\noindent \comentario{Tipo de transacciones es algo que yo pongo a mano creo, ejemplo 0:Transaccion por referidor, 1: Transferencia, 2: Canjeo de puntos}


\begin{tad}{Airmiles}



	\obs{Usuarios}{\Diccionario{id}{\tupla{MillasTotales: \Tipo{\Z}, MillasDisponiles: \Tipo{\Z}}}}

	\noindent \obs{HistorialDeTransacciones}{\seq{Transacciones}}
	\\
	\texttt{obs} Promociones : dict \{ \\
	\hspace*{1.5em} NombrePromocion: $\str$, \\
	\hspace*{1.5em} $\seq{\Tipo{FactorDePromocion},\Tipo{topeDePromocion}, \Tipo{millasRestantes}}$ \\
	\}


	\begin{proc}{ejemplo}
		{\Inout EJ: Ejemplo, \In ejemplito : \Tipo{\Z}, \In ejemplote : \Tipo{\Z}}
		{}
		\requiereLargo{ }

		\aseguraLargo{
		}
	\end{proc}
\end{tad}
\end{document}